% Archivo generado por tools/compare_suites.py: no editar a mano.
\begin{tabular}{llp{5.6cm}ccc}
\toprule
& Req. & Error introducido & \textbf{Repo} & \textbf{IA básica} & \textbf{IA estruct.} \\
\midrule
M1 & R1 & Rechaza la coordenada 3 (límite superior excluido) & \killed & \killed & \killed \\
M2 & R1 & Acepta la coordenada 0 & \killed & \survived & \killed \\
M3 & R1 & Acepta la coordenada 4 & \killed & \survived & \killed \\
M4 & R1 & Informa el eje equivocado al salir del eje X & \killed & \survived & \killed \\
M5 & R1 & Permite jugar sobre una casilla ocupada & \killed & \killed & \killed \\
M6 & R1 & El tablero nunca se considera lleno & \killed & \survived & \killed \\
M7 & R2 & El primer turno es de O & \killed & \killed & \killed \\
M8 & R2 & Los turnos no alternan & \killed & \killed & \killed \\
M9 & R2 & Una jugada inválida consume el turno & \killed & \survived & \survived \\
M10 & R3 & No reconoce la diagonal secundaria & \killed & \killed & \killed \\
M11 & R3 & No reconoce la tercera columna & \killed & \survived & \killed \\
M12 & R3 & No reconoce la fila inferior & \killed & \survived & \killed \\
M13 & R3 & Un tablero lleno es empate aunque haya línea & \killed & \survived & \survived \\
M14 & R3 & Se puede seguir jugando después de ganar & \killed & \survived & \killed \\
M15 & R3 & winner() informa un ganador cuando nadie ganó & \killed & \survived & \killed \\
M16 & R3 & reset() no reinicia el estado de la partida & \killed & \survived & \killed \\
\bottomrule
\end{tabular}
